\subsection{10.62: a periodic perfect group generated by involutions}
\label{cand:10.62}
\problemstatement{10.62}
\begin{proposition}
For every sufficiently large odd prime \(p\), there is an infinite
three-generated perfect group \(G\) of exponent \(2p\), with no subgroup
of order four, whose involutions form one conjugacy class and generate
\(G\). Any two distinct involutions have product of order \(p\).
\end{proposition}
The deep existence input is Amelio's geometric quotient construction
\cite[Proposition~6.1, Theorem~6.7, Remark~6.8]{amelio}.
We use its construction as well as its final theorem: the first quotient
may impose a prescribed eligible primitive odd-power relation.
This feature is needed to eliminate the quotient of order two.
The cone-off and small-cancellation estimates underlying that result
are imported, not reproved by the experiment.

\begin{proof}
Start with
\[
 F=\langle a,b,c\mid a^2=b^2=c^2=1\rangle=C_2*C_2*C_2
\]
acting by left multiplication on its trivalent Cayley tree with
unit-length edges. Vertex stabilizers are trivial; a nontrivial point
stabilizer is the order-two group interchanging an edge's endpoints.
A finite subgroup fixes the center of a finite invariant subtree and
therefore has order at most two.
The action is acylindrical: the number of elements moving one point
at most \(R\) is bounded by the size of a uniform vertex ball of radius
\(R+1\). This also shows that elliptic subgroups are finite.
The translations \(ab,ac\) have different endpoint pairs, so the action
is non-elementary. Every nonelliptic element translates by at least two,
and the injectivity radius is exactly two.

Amelio calls an action tame when there is no subgroup of order four
and each maximal loxodromic subgroup has finite normal subgroup of
order at most two. Here an axis stabilizer acts faithfully on a
simplicial line, since it fixes no vertex nontrivially; it is cyclic
or infinite dihedral and has trivial finite normal subgroup.
Thus the action is tame.
The additional numerical invariants in Theorem~6.7 are finite.
More precisely, its invariant \(\nu\) is at most one:
if \(h\) is loxodromic and \(\langle g,h^{-1}gh\rangle\) is elliptic,
then \(g\ne1\) would force \(h\) to fix the unique midpoint fixed by
\(g\), a contradiction. Thus \(\tau=\max\{\nu,3\}=3\).
Regard the tree as \(\delta\)-hyperbolic for
\(0<\delta<\min\{1/L_S,1/100\}\), with \(L_S\) the source's constant.
Every short element used in its overlap invariant \(\Omega\) is then
elliptic. The \(13\delta\)-enlargement of its almost-fixed set lies
within \(17\delta\) of its unique midpoint.
Distinct midpoints are at least one apart, while a non-elementary
tuple contains two distinct nontrivial elliptic elements.
Hence the relevant intersections are empty and \(\Omega=0\).
Using positive \(\delta\) respects the source's strict displacement
inequalities. Take \(Q\) to be all loxodromic elements, explicitly
permitted as a strongly stable family by Remark~6.6.

The word \(w=abc\) is primitive in its maximal loxodromic subgroup.
Its axis has periodic label pattern \(abcabc\cdots\), whose least
positive period is three; free vertex action makes its translation
stabilizer exactly \(\langle w\rangle\).
An orientation reversal would identify \(abc\) with \(cba\) up to
cyclic rotation, which is impossible. Thus the full axis stabilizer
is \(\langle w\rangle\).

Choose an initial scale \(\lambda>0\) sufficiently small that
\(3\lambda\le L_S\delta_1\), together with the initialization
conditions in Theorem~6.7. Remark~6.8 permits this smaller scale
at the cost of increasing the critical exponent.
The new injectivity radius is still \(2\lambda>0\); the required
lower bound tends to zero as the critical exponent grows.
All choices are made before choosing \(p\).
Explicitly, the source's Remark~6.8 illustrates the extra condition
\(2\lambda\le L_S\delta_1\); the word \(abc\) requires
\(3\lambda\le L_S\delta_1\) instead. The permission to decrease
\(\lambda\) is the part we use. For a sufficiently large odd prime
\(p\), use the exponent family \(\{p\}\). The primitive cyclic
subgroup \(\langle abc\rangle\) is eligible for the first family;
start with it and extend to a maximal admissible selection, choosing
one representative for each eligible maximal loxodromic subgroup
up to conjugacy. Proposition~6.1 kills \(w^p\).
Neither the extra relation nor the scale \(3\lambda\) is asserted
as part of the headline statement of Theorem~6.7.

The theorem constructs epimorphisms
\[
 F=F_0\twoheadrightarrow F_1\twoheadrightarrow F_2
 \twoheadrightarrow\cdots\twoheadrightarrow G
\]
with infinite limit, embedding the initial elliptic subgroups,
and with every initial loxodromic element mapping to an element
of order dividing \(2p\) or to an initial elliptic image.
Consequently \(a,b,c\) survive as involutions generating \(G\),
every element of \(G\) has order dividing \(2p\), and
\[
 (abc)^p=1.                                           \tag{****}
\]
Each \(F_i\) is tame and has no subgroup of order four.
This property passes to the limit here: a finite subgroup of \(G\)
has finitely many multiplication-table relations; choose lifts of
its elements, and a stage at which all those relations hold.
They define an embedding into that stage, because composition to
\(G\) is the original inclusion.
Thus \(G\) also has no subgroup of order four.

Let \(x,y\) be distinct involutions of \(G\), including any not
identified beforehand with original involutions. Put \(t=xy\).
If \(t\) had even order \(d\), then \(z=t^{d/2}\) would be an
involution centralized by \(x\).
Also \(x\notin\langle t\rangle\), since otherwise \(y\) would lie
there too and a cyclic group has only one involution.
Thus \(\langle x,z\rangle\cong C_2^2\), a contradiction.
Hence \(|t|=p\). With \(k=(p+1)/2\),
\[
 t^{-k}xt^k=xt^{2k}=y,
\]
so all involutions are conjugate. There are at least two, since
otherwise \(G\) would be cyclic of order two. The exponent is exactly
\(2p\).

In the abelianization all three generating involutions have the same
image \(\epsilon\), of order dividing two. Relation (****) gives
\(3p\epsilon=0\), and oddness implies \(\epsilon=0\).
Thus \(G\) is perfect, and in particular has no subgroup of index two.
\end{proof}

No numerical threshold for \(p\), finite presentation, or simplicity
claim follows here. The new candidate is the application with the
selected odd relation; the geometric construction is Amelio's.
Finite dihedral and abelianization controls in
\artifact{results/10.62-summary.json} check only the elementary
deductions. They do not construct this infinite group.
The complete parameter and source audit is
\artifact{research/10.62-review.md}.
